% Part: history
% Chapter: biographies
% Section: emmy-noether

\documentclass[../../../include/open-logic-section]{subfiles}

\begin{document}

\olfileid{his}{bio}{noe}

\olsection{Emmy Noether}

\olphoto{noether-emmy}{Noether}

Noether（诺特）(\textsc{ner}-ter) 于 1882 年 3 月 23 日出生在德国埃尔朗根的一个中上阶层学者家庭。尽管当时女性接受教育面临重重障碍，她仍在数学和物理学领域做出了开创性贡献，被誉为“现代代数之母”。在当时的德国，年轻女孩理应只接受艺术教育，并且不被允许进入大学预科学校。然而，在旁听了哥廷根大学和埃尔朗根大学（其父在该校担任数学教授）的课程后，随着埃尔朗根大学于 1904 年修改政策允许女生入学，Noether最终得以注册成为学生。她于 1907 年获得了数学博士学位。

尽管完全具备任职资格，Noether在职业生涯中仍遭遇了诸多阻力。1908 年至 1915 年间，她在埃尔朗根无薪教学。在此期间，她引起了当时世界顶尖数学家之一 Hilbert的注意，Hilbert 邀请她前往哥廷根。然而，女性当时被禁止获取教授职位，她只能以 Hilbert 的名义授课，再次没有薪酬。正是在这段时间里，她证明了如今所称的Noether定理，该定理至今仍在理论物理学中使用。Noether最终在 1919 年获得了任教资格。据报道，Hilbert 对大学同事持续不断的抵制回应道：“先生们，教授会不是澡堂。”

20世纪20年代后期，她专注于抽象代数的研究，她的贡献彻底改变了这一领域。在她的证明中，她经常使用所谓的升链条件，即某些集合不存在严格递增的无限链。例如，如今称为Noether环的某些代数结构具有如下性质：不存在理想的无限序列 $I_1 \subsetneq I_2 \subsetneq \dots$。这一条件可以推广到任何偏序（在代数中，它涉及由子集关系排序的理想这一特例），我们也可以考虑对偶的降链条件，即偏序中任何严格\emph{递降}的序列最终都会终止。如果一个偏序满足降链条件，我们就可以沿着该序进行归纳，类似于我们沿着~$\Nat$ 上的 $<$ 序进行归纳的方式。这样的序称为\emph{良基的}或\emph{Noether的}，相应的证明原理称为\emph{Noether归纳法}。

Noether是犹太人，1933 年纳粹掌权后，她被解除了职务。幸运的是，Noether得以移民美国，在宾夕法尼亚州的布林莫尔学院获得了一个临时职位。在此期间，她也在普林斯顿大学讲学，尽管她发现这所大学对女性并不欢迎 \citep[81]{Dick1981}。1935 年，Noether接受了切除子宫肿瘤的手术。她术后因感染去世，安葬于布林莫尔。

\begin{reading}
关于Noether的传记，可参阅 \citet{Dick1981}。圆周理论物理研究所有关于Noether生平与影响的在线讲座 \citep{Perimeter2015}。如果你读累了，《历史课上错过的事》有一期关于Noether生平与影响的播客 \citep{Frey2015}。Noether著作集的德文原版可供查阅 \citep{Noether1983}。
\end{reading}

\end{document}
